\section{Coding Models}
\sectionkicker{CODING FRONTIER}
\begin{wideflow}
{\Large\bfseries コーディングモデルの新顔――Grok 4.7}\par
\smallskip
{\color{SurveyMuted}9月21日のGrok 4.7を、値とベンチマークの非対称の留保とともに示す。}\par
\end{wideflow}

9月21日、SpaceXAIはコーディングと知識労働向けにGrok 4.7を出した\cite{grok47}。ベースモデルを大きくし、数時間かかる難題に寄せた長い強化学習を施し、Grok Botのharnessへのネイティブ対応を組み込んだとされる。Cursorの各プラン、Grok Build（無料）、APIやモデルルーター、各クラウドで利用できる。値は4.6と同じく入力100万トークン2ドル、出力6ドル、高速モードは倍速で倍額とされる。追加学習にCursorの仕事の記録（匿名化）を使ったこともモデルカードに書かれる。

ベンチマーク結果はeffort設定が揃っていない。CursorBench 4.0では46.3\%（4.7はxhigh、4.6はhigh、Solはmax、Fableはmaxの並び）とされ、DeepSWEでは71.0\%、Terminal-Bench 4.0では自社harnessで37.6\%とされる。第三者の測定では、Valsの集計が1タスクあたり費用の安さを示す一方で長時間タスクほど嵩み、Artificial Analysisの測定では同じベンチマークが26\%と自社harnessでの38\%から離れる\cite{datacampgrok47}。harness依存が見えるため、本号では優劣を言わない。

\begin{claimboundary}[資料と評価の境界]
ベースモデルの規模や強化学習の期間、harnessへの対応は、他所から確かめようのない開発元の言い分である。端末操作タスクの優劣はharnessの違いに左右されるため断定しない。公開の場の26件にはGrok 4.7の該当する投稿がなく、動向を主張する根拠には使わない。
\end{claimboundary}
